\documentclass{article}
\usepackage{xcolor}
\newcommand{\draftnote}[1]{\textcolor{red}{[Draft note: #1]}}
\newcommand{\TEKG}{TE-KG}
\newcommand{\coexpression}{co-expression}

\title{TE-KG: a provenance-aware knowledge graph and multi-layer resource for human transposable elements}
\author{Author names to be supplied\\Department and institution to be supplied\\Corresponding author details to be supplied}
\date{Working draft v0.1, 1 August 2026}

\begin{document}
\maketitle

\begin{abstract}
Human transposable element evidence is distributed across repeat classifications, reference sequences, genomic annotations, expression datasets and biomedical literature. These sources answer different questions and are difficult to examine together without obscuring their distinct units of evidence. We developed TE-KG, a human transposable element resource that connects literature-derived biomedical relations, taxonomy, representative sequence and genomic records, expression profiles and context-specific co-expression networks. In the runtime snapshot reported here, the knowledge graph contains 225 transposable element nodes, 2,308 paper nodes and 12,444 directed biological relations. Each biological relation retains a predicate and PubMed provenance. Expression data comprise 1,158 samples across normal tissue, normal primary cell and cancer cell-line contexts, and the co-expression catalogue provides approved searchable networks for 285 transposable elements and 499 genes. The web resource supports browsing, graph and path exploration, expression inspection, data download and evidence-bounded natural-language questions. TE-KG is designed to connect complementary evidence rather than collapse it: literature relations are not interpreted as causal, consensus sequences are not treated as copy-specific and co-expression edges indicate correlation rather than regulation. This structure provides traceable routes from transposable element identity to several forms of biological evidence while retaining the boundary of each source layer.
\end{abstract}

\noindent\textbf{Keywords:} transposable elements; knowledge graph; biomedical literature; expression; co-expression; database

\noindent\textbf{Database URL:} \draftnote{A stable, public, login-free production URL must be supplied before submission.}

\section{Introduction}

Transposable elements (TEs) are a major component of the human genome and are studied through several complementary representations. Repeat-family resources describe classification and consensus sequence models, genome annotations record individual or representative loci, expression datasets measure transcription in specific samples, and the biomedical literature reports associations with genes, diseases and other entities. These representations are not interchangeable. A family consensus is not the sequence of every genomic copy, an observed genomic locus does not establish activity, and a reported association does not by itself establish causality. The interpretation of TE evidence therefore depends on both the source and the unit of observation.

Established resources provide substantial depth within particular evidence types. Repbase and Dfam support repeat-family classification and reference sequence models \citep{Bao2015Repbase,Kojima2018HumanRepbase,Wheeler2013Dfam}. HERVd, dbRIP, euL1db and TranspoGene organize endogenous retrovirus sequences, retrotransposon insertion polymorphisms, sample-specific L1HS insertions or gene-context annotations, respectively \citep{Paces2002HERVd,Wang2006dbRIP,Mir2015euL1db,Levy2008TranspoGene}. More recent resources have focused on disease or expression evidence. HervD Atlas curates HERV-disease associations and exposes them through an interactive knowledge graph, CancerHERVdb concentrates on HERV expression in cancer, and TE-SCALE provides single-cell TE expression and TE-gene co-expression across human cancers \citep{Li2024HervDAtlas,Stricker2023CancerHERVdb,Meng2026TESCALE}. These specialist resources answer important questions at resolutions that a general integration resource should not claim to replace.

The remaining practical problem is cross-layer navigation. A researcher starting from a TE name may need to establish its classification, inspect a reference sequence or genomic record, compare expression contexts, identify correlated genes and open the papers supporting a disease or gene relation. When these steps require separate resources, names and evidence boundaries must be reconciled repeatedly. Conversely, merging heterogeneous records into a single undifferentiated graph risks making weak or indirect evidence appear stronger than it is. The useful integration task is therefore not only to connect records, but also to retain their provenance and interpretation limits.

Here we describe TE-KG, a human TE-centred resource that connects literature-derived biomedical relations, taxonomy, representative genomic and sequence records, expression profiles and context-specific co-expression networks. TE-KG provides graph, path, table, visualization, download and natural-language access routes to these layers. The central design principle is that integration should make evidence jointly accessible without erasing what each source can support. We report the current database content, data-processing and runtime architecture, and the principal scientific access workflows. We also state the limitations that constrain biological interpretation and the additional release information required before public deployment.


\section{Materials and methods}

\subsection{Scope and evidence model}

TE-KG was designed around human TEs as shared identifiers connecting otherwise distinct evidence layers. The literature layer stores biomedical entities and relations extracted from publications. The taxonomy layer represents classification relationships among TE records. Sequence and genomic records provide reference or representative information derived from different upstream sources. Expression matrices retain sample-context measurements, and co-expression networks summarize context-specific rank correlations between TE and gene features. The implementation does not assign a common evidential meaning to these layers. Instead, each runtime route identifies the source type and restricts the language used to describe it.

\subsection{Literature collection, extraction and normalization}

PubMed was queried for human TE literature using a recorded combination of TE names, aliases and TE-related terms. The archived workflow covers publications available through 13 April 2026. Candidate records were screened against a TE whitelist and normalized before inclusion. The final retained corpus comprised 2,308 papers, matching the number of Paper nodes in the current Neo4j database snapshot.

The retained text and metadata were processed with constrained DeepSeek-V3 prompts to identify TE-centred biomedical statements and their participating entities. Extracted entities were mapped to 11 biomedical categories: carbohydrate, disease, function, gene, lipid, mutation, peptide, pharmaceutical, protein, RNA and toxin. Relation records retained a normalized predicate, one or more supporting PubMed identifiers (PMIDs) and the number of supporting PMIDs. Automated normalization included name matching using a Ratcliff--Obershelp similarity threshold of 80\%, followed by manual mapping where automated resolution was insufficient. Manual curation was used to review retained entities and relations before graph import.

The archived method record reports three manual audits of 50 records each. Those percentages are not used as a quantitative performance claim in this manuscript draft because the sampled identifiers, sampling seed, exact denominators and intermediate screening arithmetic have not yet been recovered consistently. \draftnote{Recover the literature query string, exact DeepSeek-V3 endpoint/model version, intermediate record counts and sampled audit decisions before locking this subsection.}

\subsection{TE classification, sequences and genomic records}

TE identity and classification records were assembled from a RepBase-derived local source, whereas genomic positions were obtained from a RepeatMasker-derived (RMSK) annotation source. Repbase is an established reference for repetitive-element families and consensus sequences \citep{Bao2015Repbase,Kojima2018HumanRepbase}; Dfam provides a complementary family-model framework \citep{Wheeler2013Dfam}. TE-KG preserves the distinction between these source roles. RepBase-derived sequences are described as consensus or reference sequences, and RMSK-derived positions are described as representative genomic records. Neither is presented as an exhaustive inventory of active, polymorphic or sample-specific insertions.

Taxonomy is served from the Neo4j-backed API rather than from a second browser-side truth source. Classification edges use the explicit graph relations \texttt{CLASSIFIED\_AS}, \texttt{SUBFAMILY\_OF} and \texttt{HAS\_SUBCATEGORY}. The taxonomy API can expose the current tree for all supported TE records or the RMSK plus RepBase view used by the web interface. \draftnote{Insert the exact upstream RepeatMasker build, RepBase release, acquisition dates and licence/redistribution terms.}

\subsection{Expression datasets and context definitions}

The expression layer comprises three active matrices. The normal-tissue matrix contains 37,868 features and 205 samples; the normal-primary-cell matrix contains 37,868 features and 307 samples; and the cancer-cell-line matrix contains 37,868 features and 646 samples. Together, the matrices contain 1,158 samples. Feature identifiers include TE or repeat features and genes, so the matrices are not described as gene-only expression data.

The normal-tissue data are associated with the E-MTAB-1733 and E-MTAB-2836 studies \citep{Edqvist2015SkinTranscriptome,Uhlen2015HumanProteome}. Cancer-cell-line data are associated with PRJNA523380 and the Cancer Cell Line Encyclopedia \citep{Ghandi2019CCLE}. The current normal-primary-cell matrix is associated with SRP013565, an ENCODE RNA-seq study that contains heterogeneous experiments. \draftnote{Freeze an accession-to-local-sample manifest for all three matrices and identify the precise SRP013565 subset and primary citation.}

Metadata auditing identified duplicated run identifiers in the normal-tissue metadata and five matrix runs without matching metadata (ERR579126, ERR579137, ERR579144, ERR579145 and ERR579154). These records remain visible as a limitation and must be resolved or explicitly excluded in the public release manifest. The runtime key \texttt{normal\_cell\_line} refers to the normal-primary-cell context; manuscript labels use \emph{normal primary cell} to avoid confusing it with transformed cell lines.

\subsection{Context-specific co-expression networks}

Co-expression was computed separately within normal tissue, normal primary cell and cancer cell line contexts. Abundance values were transformed as $\log_{2}(\mathrm{count}+1)$, and TE-gene associations were measured using Spearman rank correlation. Candidate edges were retained when $|r| \geq 0.4$ and the Benjamini--Hochberg false discovery rate (FDR) was at most 0.05 \citep{Benjamini1995FDR}. Positive retained edges were used for Louvain community detection with random seed 42 and resolution 1.8 \citep{Blondel2008Louvain}. The resource presents these edges as statistical correlations, not as evidence that a TE regulates a gene.

Full offline network outputs are retained as provenance and importer inputs. The production interface serves approved MySQL-backed display networks selected from those results. The current display contract limits a returned network to at most 50 nodes and 150 edges, allowing responsive inspection without describing the displayed subgraph as the complete correlation network. The catalogue version reported in this manuscript is \texttt{v1\_abs0.4\_fdr0.05\_res1.8}.

\subsection{Graph, relational storage and web application}

The production graph is stored in Neo4j database \texttt{tekg3}. MySQL stores the versioned Browse catalogue and the expression and approved co-expression datasets. PHP APIs mediate access to both database systems, and browser JavaScript renders tables, charts, the taxonomy views and interactive G6 or Canvas graphs. This separation prevents large expression matrices and display-network records from being forced into the literature graph while retaining a common TE-centred interface.

The knowledge-graph interface supports entity search, relation filtering, interactive node expansion, relation legends and PMID-linked evidence. The Path interface searches for connections between selected entities. Browse, taxonomy and Expression provide structured access to TE identity and quantitative profiles, whereas Download exposes reusable graph, taxonomy and expression artifacts. All manuscript-facing counts were obtained from dated read-only database queries or live API responses rather than copied from interface labels.

\subsection{Evidence-bounded natural-language retrieval}

Agent and DeepThink provide two natural-language access modes over the local resource. Both use a catalogue of 12 evidence-retrieval roles covering entity resolution, graph, path, sequence, taxonomy, genome, expression, co-expression and literature tasks. The language model selects evidence according to the question, while the server enforces route restrictions, evidence sufficiency checks and bounded fallback behavior. Answers can include PubMed links and can use previous turns from the current browser conversation. Conversation context is not persisted across a reload, new tab or new session.

The intelligent interface is treated as an access layer rather than a primary scientific result. Internal routing labels, plugin names and developer diagnostics are removed from user-facing prose. The current manuscript therefore reports implemented behavior and test coverage but does not claim a single overall accuracy value. Targeted limitations, including title-level literature retrieval in some cases and incomplete disease-qualified searches, remain disclosed.

\subsection{Verification and snapshot policy}

Database counts were verified on 31 July 2026 using read-only Neo4j queries, live PHP APIs and repository checks. A biological relation was counted once in its stored direction. The alternative undirected traversal count was rejected because it counted each stored directed relation twice. API contracts, runtime database configuration, graph rendering, taxonomy truth, co-expression catalogue parity and Agent/DeepThink behavior were checked with the repository's automated and browser-oriented regression scripts. Environment-specific response times were not used as biological or performance claims.


\section{Results}

\subsection{A literature-linked human TE graph with explicit provenance}

The current Neo4j snapshot contains 225 TE nodes and 2,308 Paper nodes. It also contains 12,444 directed \texttt{BIO\_RELATION} relationships connecting TEs to biomedical entities. Every current biological relation has a normalized predicate, at least one PMID field and a supporting-PMID count. These properties allow a displayed edge to be traced to the literature record used to support it, although provenance does not convert an extracted association into causal evidence.

The biomedical schema contains 11 entity categories. The most numerous current labels are Function (3,683 nodes), Gene (1,280), Protein (1,089), Disease (676), RNA (588), Mutation (377) and Pharmaceutical (293). Smaller categories include Toxin (67), Lipid (26), Peptide (23) and Carbohydrate (12). DiseaseCategory (767), Paper (2,308) and taxonomy-related nodes are represented separately because they serve organizational or provenance roles rather than the 11 biomedical entity categories.

\begin{table*}[!t]
\caption{Current TE-KG content snapshot. Counts describe different runtime units and must not be summed. The snapshot was verified on 31 July 2026.\label{tab:snapshot}}
\begin{tabularx}{\textwidth}{@{}l r X X@{}}
\toprule
Component and unit & Count & Definition & Source and boundary \\
\midrule
Neo4j TE nodes & 225 & TE entities in the \texttt{tekg3} graph & Live label-count query; not identical to Browse entries \\
Neo4j Paper nodes & 2,308 & Retained publications represented in the graph & Live label-count query; agrees with the final literature corpus \\
Directed biological relations & 12,444 & Stored \texttt{BIO\_RELATION} relationships & Counted once in the stored direction; all current records contain predicate and PMID provenance fields \\
Biomedical entity categories & 11 & Carbohydrate, Disease, Function, Gene, Lipid, Mutation, Peptide, Pharmaceutical, Protein, RNA and Toxin & Excludes Paper, DiseaseCategory and taxonomy labels \\
Browse catalogue entries & 276 & Versioned TE catalogue records & MySQL-backed Browse API; a different record unit from graph TE nodes \\
Taxonomy-classified TEs & 192 of 225 & TE nodes with a taxonomy class in the live summary & Neo4j-backed taxonomy API \\
Expression samples & 1,158 & 205 normal tissue, 307 normal primary cell and 646 cancer cell line & Three source groups; not a matched cross-context cohort \\
Searchable co-expression entries & 784 & 285 TE and 499 Gene entries & Approved display catalogue across three contexts; not the number of full offline-network features \\
Download files & 10 & Six expression, two graph and two taxonomy files & Current web route; public archival release still required \\
\botrule
\end{tabularx}
\end{table*}



\subsection{Classification and catalogue views retain distinct record units}

The live taxonomy summary contains 225 TEs, of which 192 have an assigned taxonomy class. The summary also reports 188 standard leaves and separates a 36-record all-source tree from a 189-record RMSK plus RepBase tree. These values describe API-defined taxonomy views rather than independent biological totals. The Browse catalogue contains 276 versioned entries. This larger Browse count is reported separately from the Neo4j TE-node count because the two components expose different record units and should not be combined into a single TE total.

The classification interface offers both a hierarchical tree and a force-directed graph while using the same API-backed taxonomy data. Its legend can temporarily hide or restore TE classes without changing the underlying result set. This interaction supports inspection of dense classifications, but it does not introduce a second taxonomy source or modify classification membership.

\subsection{Expression and co-expression add three context classes}

The active expression matrices contain 205 normal-tissue, 307 normal-primary-cell and 646 cancer-cell-line samples. These 1,158 samples provide three context classes for TE and gene expression inspection. The resource keeps the contexts separate because they originate from different studies and sampling designs. The current data do not support treating them as a matched experimental comparison.

The production co-expression catalogue spans the same three context classes and includes approved searchable entries for 285 TEs and 499 genes. Across contexts, 849 display recommendations are available: 17 are marked as core cases, 287 as high-confidence entries, 542 as searchable entries and three as not recommended for default display. These tiers govern presentation and searchability; they are not statistical significance categories. All displayed edges already satisfy the fixed correlation and FDR thresholds, and a returned display network is a bounded view of a larger offline result.

\subsection{Connected interfaces expose complementary evidence routes}

TE-KG provides several entry routes suited to different questions. Browse and Search resolve TE names and lead to structured identity records. Knowledge Graph displays literature-linked biomedical relations and exposes PubMed evidence. Path searches for a connection between selected entities. Expression shows context-specific abundance profiles, and Co-expression displays approved TE-gene correlation neighbourhoods. Download currently exposes ten files comprising six expression artifacts, two graph artifacts and two taxonomy artifacts.

These routes are connected at the level of TE identity but retain separate evidence semantics. For example, opening a disease-linked graph edge gives access to the supporting paper rather than a claim that the TE causes the disease. Likewise, opening a co-expression neighbourhood reports the context and correlation threshold rather than a regulatory mechanism. This separation is part of the user-visible design, not only an implementation detail.

\subsection{Natural-language follow-up remains bounded to the current session}

Agent and DeepThink can retrieve several local evidence layers in response to an English question and can interpret a follow-up turn within the current conversation. PubMed identifiers in generated answers are rendered as links to PubMed. The writing layer translates evidence limitations into ordinary scientific language instead of displaying internal routing, confidence labels or other developer-facing diagnostics. For example, an answer can describe a relation as an association that requires experimental validation without asking the reader to interpret an internal status code.

The current evaluation supports functional coverage but not a global accuracy claim. A fixed 13-question baseline exercised all 12 retrieval roles, and subsequent regression checks protected graph-routing, literature-link and multi-plugin behavior after targeted fixes. A larger historical 36-question run included failures that were later corrected, but the entire suite has not been rerun as one version-locked benchmark. We therefore report this interface as a bounded retrieval route and retain individual known limitations rather than publishing the historical pass fraction as current performance.

\subsection{Cross-layer L1HS use case}

L1HS was selected as the principal worked example because it can be resolved across the current classification, graph, sequence, genome, expression and co-expression interfaces. The final figure will preserve the distinct result type in each panel and link all literature statements to verified PMIDs. \draftnote{Freeze the exact L1HS records, expression source table, co-expression neighbourhood and one or two PMID-supported graph relations before this subsection is considered manuscript-ready. Do not fill missing layers from general model knowledge.}

\section{Discussion}

TE-KG addresses a practical integration problem in human TE research: evidence about the same named element is distributed across family references, genomic annotations, expression matrices and biomedical papers. The resource links these layers through shared TE identities and connected access workflows. Its main contribution is not the use of a graph in isolation. Rather, it is the combination of a PMID-bearing literature graph with classification, representative sequence and genomic records, three expression contexts and context-specific co-expression views, while keeping the interpretation boundary of each layer visible.

This positioning is deliberately complementary to specialist resources. Repbase and Dfam provide substantially deeper repeat-family references and models \citep{Bao2015Repbase,Wheeler2013Dfam}. dbRIP and euL1db address insertion polymorphisms and sample-specific L1HS records at a resolution not provided by TE-KG \citep{Wang2006dbRIP,Mir2015euL1db}. HervD Atlas offers manually curated depth for HERV-disease associations, and TE-SCALE provides single-cell cancer expression and co-expression at a scale beyond the bulk datasets described here \citep{Li2024HervDAtlas,Meng2026TESCALE}. TE-KG is therefore most useful when a question crosses evidence types, not when a user needs the deepest available record within one specialist domain.

The literature graph offers a clear provenance advantage and a clear limitation. Retaining predicates and PMIDs makes graph edges inspectable and allows users to return to the source publication. However, the current corpus was produced through a combination of automated screening, constrained language-model extraction, normalization and manual curation. Errors can arise from retrieval, abbreviation resolution, entity normalization or interpretation of the source text. PMID presence establishes traceability, not correctness or causality. A reproducible manual audit with preserved sample identifiers is required before the extraction pipeline can support a quantitative accuracy statement.

Expression and co-expression require similar caution. TE expression estimates are sensitive to repetitive sequence and quantification choices \citep{Lanciano2020TEExpression}. The three TE-KG contexts originate from different studies and cannot be interpreted as a matched normal-to-cancer experiment. The normal-tissue metadata also contain duplicated identifiers and five unmatched matrix runs. Co-expression further reduces the data to thresholded pairwise correlations and modules. FDR control limits the expected proportion of false discoveries among retained statistical tests, but it does not establish a regulatory interaction. The display catalogue adds another selection layer because it presents bounded, approved subgraphs rather than every offline edge.

The intelligent interface lowers the effort needed to retrieve evidence across several stores, but it should remain secondary to source-visible database access. Session-scoped follow-up makes short research conversations possible without building a persistent user profile. At the same time, answer quality depends on entity resolution, plugin selection, the availability of local evidence and the language model's synthesis. The interface can omit relevant literature or produce an incomplete answer even when the underlying resource contains useful records. Keeping PubMed links visible and translating uncertainty into reader-facing language are therefore more important than presenting internal reasoning steps.

Several requirements remain before submission and public release. The exact RepeatMasker and RepBase releases, acquisition dates and redistribution terms must be recorded. The expression accession-to-sample manifest must be frozen, including the SRP013565 subset and the five normal-tissue metadata mismatches. A stable public URL, versioned code and data deposits, licences and an update policy are also needed. The current \emph{Database} instructions require a free, login-free resource that remains available for at least two years, so a local deployment and Download page alone are insufficient.

The most informative next evaluation is a reproducible cross-layer use case rather than a larger catalogue of interface screenshots. A locked L1HS example can test whether a user can move from identity to classification, representative sequence or location, expression, co-expression and literature evidence without losing provenance or overstating the result. A second case involving a HERV or DNA transposon would help determine whether the workflow generalizes beyond the best represented LINE example. These cases should report missing evidence as part of the result rather than supplementing it with unverified generated text.

\section{Conclusion}

TE-KG provides a provenance-aware route through several forms of human TE evidence. Its current knowledge graph links 225 TE nodes to 2,308 papers through 12,444 directed, PMID-bearing biological relations, while separate taxonomy, sequence, genomic, expression and co-expression stores support complementary questions. Connected web and natural-language workflows make these layers accessible without assigning them a common evidential meaning. The resource is therefore best understood as an integration and navigation layer: literature relations remain associations, reference sequences remain distinct from individual genomic copies and co-expression remains correlation. With the remaining provenance, release and worked-use-case requirements resolved, this structure can support traceable investigation of human TEs across data types while directing users back to the underlying records.


\section{Data availability}

The current development version exposes graph, taxonomy and expression downloads through the TE-KG web interface. \draftnote{Replace this working statement with a stable public Database URL, versioned repository and archival deposit identifiers, file-level licences, release version, update policy and accession-to-sample manifests. The final resource must be freely accessible without login or registration.}

\section{Supplementary data}

\draftnote{Supply the frozen source tables, query files, data manifests, literature-audit records, co-expression sensitivity outputs and detailed verification results described in the table plan.}

\section{Funding}

\draftnote{Author input required, including funder names, grant numbers and recipient initials.}

\section{Acknowledgements}

\draftnote{Author input required.}

\section{Conflict of interest}

\draftnote{Author input required.}

% Author Contributions is deliberately omitted at the author's request and
% must not be inferred from repository history.


\bibliography{../references}
\end{document}

